% TOP_PROBLEM: {"id": 8400010, "title": "O9 — Counting distinct prime factors", "category_id": 3, "category": {"id": 3, "name": "graph_theory", "display_name": "Graph Theory", "description": "Problems involving graphs, networks, and their properties.", "slug": "graph-theory", "order_index": 3, "created_at": "2026-07-31T15:26:25.671Z"}, "status": "open"}
% FIRST_POSTED: null
% ENTRY_KIND: statement_only
% Imported from: batch03.tex; UnsolvedMath ID 8400010
\documentclass[11pt]{article}
\usepackage[a4paper,margin=27mm]{geometry}
\usepackage{amsmath,amssymb,mathtools}
\usepackage{fontspec,unicode-math}
\setmainfont{Latin Modern Roman}
\setmathfont{Latin Modern Math}
\usepackage{xurl,hyperref}
\hypersetup{colorlinks=true,urlcolor=blue,linkcolor=blue}
\newcommand{\sref}[2]{\hyperlink{source:#1:#2}{[S#2]}}
\newcommand{\cataloguescope}[1]{\paragraph{Mathematical question.} #1}
\begin{document}
% UnsolvedMath ID 8400010; source catalogue code AMR-083-0010
\setcounter{section}{8400009}
\section{O9 — Counting distinct prime factors}\label{8400010}
{\small\sffamily UnsolvedMath ID 8400010 · Graph Theory}
\subsection{Definitions and mathematical statement}

\paragraph{Shared notation.}$\mathbb N=\{1,2,\ldots\}$, and $\mathbb Z,\mathbb Q,\mathbb R,\mathbb C$ denote the integers, rationals, reals and complex numbers. Any other conventions are specified locally.


Integers are supplied in binary; polynomial time means a polynomial number of classical bit operations in total input length. Define $\omega(n)=|\{p:p\text{ is prime and }p\mid n\}|$ for $n\ge1$, with $\omega(1)=0$.
\paragraph{Statement recorded in the supplied source.}
Is there a deterministic polynomial-time algorithm which outputs $\omega(n)$ for every positive integer $n$? \sref{8400010}{2}
\subsection{Short English statement}
Can the number of distinct prime divisors be computed deterministically in time polynomial in the binary input length?
\subsection{Sources}
\begin{description}
\item[\hypertarget{source:8400010:1}{S1}] ULAM.AI and the UnsolvedMath Contributors, UnsolvedMath: A Curated Collection of Open Mathematics Problems, record 8400010 (AMR-083-0010). CC BY 4.0. \url{https://huggingface.co/datasets/ulamai/UnsolvedMath}.
\item[\hypertarget{source:8400010:2}{S2}] Leonard M. Adleman and Kevin S. McCurley, Open Problems in Number Theory Complexity, II, in Algorithmic Number Theory, LNCS 877 (1994), pp. 291--322, C9/O9, manuscript p. 9 (LNCS p. 299). \url{https://mccurley.org/papers/open.ps.gz}.
\end{description}
\label{end:8400010}
\end{document}
